\section{Part II.b: Labels and Charges}

\subsection{Gauge}

Internal symmetry in this framework has no independent space to act on. The only structure at each point is the directional scale pattern. Consequently, any symmetry that is not a motion of the substrate must be a relabelling of directions that leaves every local unit unchanged. This module develops the residual symmetries: they form a group determined entirely by the scale pattern, with no freedom to choose a gauge group. The key result is that two directions lie in the same multiplet orbit precisely when they carry the same scale.

\begin{definition}[\lean{Gauge.stabilizer}]
Let $s : \text{Fin } n \to A^\times$ be a scale pattern on $n$ directions, where $A$ is a commutative ring. The \emph{stabilizer} $\mathrm{stabilizer}(s)$ is the subgroup of permutations $\pi \in \mathrm{Equiv.Perm}(\text{Fin } n)$ satisfying
\[ \forall a \in \text{Fin } n, \quad s(\pi a) = s(a). \]
These are the relabellings of directions that leave every local unit unchanged.
\end{definition}

\begin{theorem}[\lean{Gauge.mem\_stabilizer\_iff}]
A permutation $\pi$ belongs to $\mathrm{stabilizer}(s)$ if and only if it fixes every scale value: $\pi \in \mathrm{stabilizer}(s) \iff \forall a, \, s(\pi a) = s(a)$.
\end{theorem}

\begin{proof}
Immediate from the definition.
\end{proof}

\begin{theorem}[\lean{Gauge.stabilizer\_eq\_bot\_of\_injective}]
If the scale pattern $s : \text{Fin } n \to A^\times$ is injective (all directions carry different units), then $\mathrm{stabilizer}(s) = \{1\}$ (the trivial group).
\end{theorem}

\begin{proof}
Suppose $\pi \in \mathrm{stabilizer}(s)$. Then for all $a \in \text{Fin } n$, we have $s(\pi a) = s(a)$. Since $s$ is injective, this implies $\pi a = a$ for all $a$, so $\pi = 1$. Conversely, the identity always belongs to the stabilizer.
\end{proof}

\begin{theorem}[\lean{Gauge.stabilizer\_eq\_top\_of\_isotropic}]
If the scale pattern $s : \text{Fin } n \to A^\times$ is isotropic (all directions carry the same unit: $\forall a, b, \, s a = s b$), then $\mathrm{stabilizer}(s)$ is the full permutation group.
\end{theorem}

\begin{proof}
For any permutation $\pi$, and any direction $a$, we have $s(\pi a) = s(a)$ by isotropy. Thus $\pi \in \mathrm{stabilizer}(s)$ for all $\pi$.
\end{proof}

\begin{definition}[\lean{Gauge.SameOrbit}]
Two directions $a, b \in \text{Fin } n$ lie in the \emph{same orbit} under the scale pattern $s$, written $\mathrm{SameOrbit}(s, a, b)$, if there exists a permutation $\pi \in \mathrm{stabilizer}(s)$ such that $\pi a = b$.
\end{definition}

\begin{theorem}[\lean{Gauge.sameOrbit\_iff\_eq\_scale}]
Assume decidable equality on $\text{Fin } n$. Two directions $a, b$ lie in the same orbit under $s$ if and only if they carry the same local unit:
\[ \mathrm{SameOrbit}(s, a, b) \iff s(a) = s(b). \]
Consequently, multiplets are exactly the degeneracy blocks of the scale pattern, and multiplet sizes are the sizes of those blocks.
\end{theorem}

\begin{proof}
Forward direction: If $\pi a = b$ for some $\pi \in \mathrm{stabilizer}(s)$, then $s(b) = s(\pi a) = s(a)$ by the definition of the stabilizer.

Reverse direction: If $s(a) = s(b)$, consider the transposition $\pi = \mathrm{swap}(a, b)$ that exchanges $a$ and $b$ while fixing all other elements. We claim $\pi \in \mathrm{stabilizer}(s)$. For any direction $c$:
\begin{itemize}
\item If $c = a$, then $s(\pi c) = s(b) = s(a) = s(c)$.
\item If $c = b$, then $s(\pi c) = s(a) = s(b) = s(c)$.
\item If $c \neq a$ and $c \neq b$, then $\pi c = c$, so $s(\pi c) = s(c)$.
\end{itemize}
Thus $\pi$ stabilizes $s$, and $\pi a = b$.
\end{proof}

\begin{theorem}[\lean{Gauge.internal\_symmetry\_determined}]
Assume decidable equality on $\text{Fin } n$. If two scale patterns $s, t : \text{Fin } n \to A^\times$ are identical in terms of which directions carry the same unit (i.e., $\forall a, b, \, (s a = s b) \iff (t a = t b)$), then two directions $a, b$ lie in the same orbit under $s$ if and only if they lie in the same orbit under $t$.
\end{theorem}

\begin{proof}
By \lean{Gauge.sameOrbit\_iff\_eq\_scale}, $\mathrm{SameOrbit}(s, a, b) \iff s(a) = s(b)$ and $\mathrm{SameOrbit}(t, a, b) \iff t(a) = t(b)$. The hypothesis ensures these two conditions are equivalent, so the orbit structures are identical.
\end{proof}

\subsection{Transport}

In this framework, a connection is not an independent field: it is determined by the task of matching scales from one point to another. A transport from scale pattern $s$ to $t$ is a relabelling that makes the two patterns agree when read through the same directional index. When the scale pattern is injective (all directions different), the transport is unique. When there is degeneracy, transports form a coset of the stabilizer — exactly the gauge freedom created by that degeneracy.

\begin{definition}[\lean{Transport.Transports}]
A permutation $\pi \in \mathrm{Equiv.Perm}(\text{Fin } n)$ \emph{transports} the pattern $s$ to the pattern $t$ if the relabelled pattern matches:
\[ \forall a \in \text{Fin } n, \quad t(\pi a) = s(a). \]
The permutation $\pi$ is a valid connection between the two scale fields.
\end{definition}

\begin{theorem}[\lean{Transport.transports\_self\_iff}]
Transporting a scale pattern $s$ to itself is exactly being in its stabilizer:
\[ \mathrm{Transports}(s, s, \pi) \iff \pi \in \mathrm{stabilizer}(s). \]
\end{theorem}

\begin{proof}
Both conditions state $\forall a, \, s(\pi a) = s(a)$.
\end{proof}

\begin{theorem}[\lean{Transport.transport\_unique\_of\_nondegenerate}]
If $t$ is injective and two permutations $\pi, \tau$ both transport $s$ to $t$, then $\pi = \tau$. That is, the connection is a function of the scale field alone when there is no degeneracy.
\end{theorem}

\begin{proof}
Suppose $t(\pi a) = s(a)$ and $t(\tau a) = s(a)$ for all $a$. Then $t(\pi a) = t(\tau a)$ for all $a$. Since $t$ is injective, $\pi a = \tau a$ for all $a$, so $\pi = \tau$.
\end{proof}

\begin{theorem}[\lean{Transport.transport\_coset}]
If $\pi$ and $\tau$ both transport $s$ to $t$, then $\pi \tau^{-1} \in \mathrm{stabilizer}(t)$. That is, two transports differ by an element of the stabilizer, and only by that.
\end{theorem}

\begin{proof}
For any direction $b$, we have
\[ t(\pi(\tau^{-1} b)) = s(\tau^{-1} b). \]
But $\tau$ transports $s$ to $t$, so $\tau(\tau^{-1} b) = b$ and $s(\tau^{-1} b) = t(b)$. Therefore $t(\pi(\tau^{-1} b)) = t(b)$, confirming that $\pi \tau^{-1} \in \mathrm{stabilizer}(t)$.
\end{proof}

\begin{theorem}[\lean{Transport.transports\_of\_stabilizer}]
If $\pi$ transports $s$ to $t$ and $\rho \in \mathrm{stabilizer}(t)$, then $\rho \pi$ also transports $s$ to $t$.
\end{theorem}

\begin{proof}
For any direction $a$,
\[ t(\rho(\pi a)) = t(\pi a) = s(a), \]
where the first equality uses $\rho \in \mathrm{stabilizer}(t)$ and the second uses the fact that $\pi$ transports $s$ to $t$.
\end{proof}

\begin{theorem}[\lean{Transport.transport\_unique\_iff\_stabilizer\_trivial}]
The transport from any pattern $s$ to $t$ is unique if and only if $\mathrm{stabilizer}(t) = \{1\}$:
\[ \left(\forall s, \pi, \tau, \, \mathrm{Transports}(s, t, \pi) \land \mathrm{Transports}(s, t, \tau) \implies \pi = \tau\right) \iff \mathrm{stabilizer}(t) = \{1\}. \]
\end{theorem}

\begin{proof}
Forward direction: Suppose transport is unique. Then $\mathrm{Transports}(t, t, \rho)$ and $\mathrm{Transports}(t, t, 1)$ both hold (the first by assumption, the second by reflexivity), so $\rho = 1$. Thus $\mathrm{stabilizer}(t) = \{1\}$.

Reverse direction: Assume $\mathrm{stabilizer}(t) = \{1\}$. If $\pi$ and $\tau$ both transport $s$ to $t$, then by the previous theorem, $\pi \tau^{-1} \in \mathrm{stabilizer}(t) = \{1\}$, so $\pi \tau^{-1} = 1$ and $\pi = \tau$.
\end{proof}

\begin{theorem}[\lean{Transport.nondegenerate\_no\_gauge\_freedom}]
If $t$ is injective, then $\mathrm{stabilizer}(t) = \{1\}$ and for any pattern $s$, the transport from $s$ to $t$ is unique.
\end{theorem}

\begin{proof}
Follows directly from \lean{Transport.transport\_unique\_of\_nondegenerate} and \lean{Gauge.stabilizer\_eq\_bot\_of\_injective}.
\end{proof}

\subsection{Axis}

A scale is a ratio of lengths. This simple fact has profound consequences: a scale reading cannot distinguish a direction from its reverse. Therefore, the order parameter is not a vector but a *projective* object — an unoriented axis. Projective order parameters carry both point defects and line defects; the framework thus predicts point particles, not just strings. This module establishes the projectivity and characterizes which scale patterns support a well-defined order parameter.

\begin{theorem}[\lean{Axis.bareForm\_neg}]
The bare scale form $\mathrm{bareForm}(\eta, v) = \sum_i \eta_i v_i^2$ is even in the direction: reversing $v$ to $-v$ leaves it unchanged.
\end{theorem}

\begin{proof}
\[ \mathrm{bareForm}(\eta, -v) = \sum_i \eta_i (-v_i)^2 = \sum_i \eta_i v_i^2 = \mathrm{bareForm}(\eta, v). \]
\end{proof}

\begin{theorem}[\lean{Axis.physForm\_neg}]
The measured scale form $\mathrm{physForm}(s, \eta, v)$ (which incorporates the scale field $s$) is also even in the direction.
\end{theorem}

\begin{proof}
The measured form is derived from the bare form, so it inherits the evenness property.
\end{proof}

\begin{definition}[\lean{Axis.Indistinguishable}]
Two directions $v, w \in \text{Fin } n \to A$ are \emph{scale-indistinguishable} if every scale measurement gives the same reading in any signature:
\[ \forall \eta \in \text{Fin } n \to A, \quad \mathrm{bareForm}(\eta, v) = \mathrm{bareForm}(\eta, w). \]
\end{definition}

\begin{theorem}[\lean{Axis.neg\_indistinguishable}]
A direction and its reverse are indistinguishable: $\mathrm{Indistinguishable}(v, -v)$ for all $v$.
\end{theorem}

\begin{proof}
By \lean{Axis.bareForm\_neg}, $\mathrm{bareForm}(\eta, -v) = \mathrm{bareForm}(\eta, v)$ for all signatures $\eta$.
\end{proof}

\begin{theorem}[\lean{Axis.physForm\_indistinguishable}]
The projectivity is preserved by the measured form: $\mathrm{physForm}(s, \eta, v) = \mathrm{physForm}(s, \eta, -v)$ for all scales $s$ and signatures $\eta$.
\end{theorem}

\begin{proof}
Since the measured form is derived from the bare form and the bare form is even, the measured form inherits evenness.
\end{proof}

\begin{theorem}[\lean{Axis.isotropic\_no\_order\_parameter}]
If the scale pattern is isotropic ($\forall a, b, \, s(a) = s(b)$), then the stabilizer is the full permutation group, so there is no order parameter and no defects.
\end{theorem}

\begin{proof}
This follows directly from \lean{Gauge.stabilizer\_eq\_top\_of\_isotropic}. When every direction carries the same scale, every relabelling is a symmetry, leaving no variation to be mapped across space.
\end{proof}

\begin{theorem}[\lean{Axis.nondegenerate\_full\_frame}]
If the scale pattern is injective (all directions carry different scales), then the stabilizer is trivial, and the order parameter is the full frame, not a single axis.
\end{theorem}

\begin{proof}
This follows from \lean{Gauge.stabilizer\_eq\_bot\_of\_injective}. A fully non-degenerate pattern leaves no directions to permute.
\end{proof}

\begin{theorem}[\lean{Axis.uniaxial\_stabilizer}]
Assume decidable equality on $\text{Fin } n$. Suppose the scale pattern $s$ has the form: one special direction $t$ carries a unique scale, while all other directions agree:
\[ \forall a, b \neq t, \quad s(a) = s(b). \]
Then two non-special directions $a, b \neq t$ lie in the same orbit, so the order parameter is a single unoriented axis.
\end{theorem}

\begin{proof}
By \lean{Gauge.sameOrbit\_iff\_eq\_scale}, two directions $a, b$ lie in the same orbit iff $s(a) = s(b)$. The hypothesis gives $s(a) = s(b)$ whenever $a, b \neq t$, so all non-special directions form one orbit.
\end{proof}

\begin{theorem}[\lean{Axis.uniaxial\_axis\_distinguished}]
Assume decidable equality on $\text{Fin } n$. In a uniaxial configuration, the distinguished direction $t$ is not in the same orbit as any other direction: if $s(a) \neq s(t)$, then $\neg \mathrm{SameOrbit}(s, a, t)$.
\end{theorem}

\begin{proof}
By \lean{Gauge.sameOrbit\_iff\_eq\_scale}, $\mathrm{SameOrbit}(s, a, t)$ would imply $s(a) = s(t)$, contradicting the hypothesis.
\end{proof}

\begin{theorem}[\lean{Axis.scalar\_picture\_is\_the\_isotropic\_locus}]
The scalar picture (where a single number represents the scale) and the directional picture (where each direction has its own scale) disagree on defect classification. The scalar picture corresponds exactly to the isotropic locus, and the confusion noted in the introduction is the same one already diagnosed for geometry: treating an isotropic special case as the general framework.
\end{theorem}

\begin{proof}
By \lean{Axis.isotropic\_no\_order\_parameter}, full isotropy gives stabilizer $= $ full permutation group, so no order parameter and no defects. This is the scalar picture's limiting case, not the generic prediction of the directional framework.
\end{proof}

\subsection{Defect}

Particles emerge from scale defects: topological failure of the log-scale to close up around a loop when measured modulo a discrete period $\Delta$. The winding number is an integer and is topologically stable. Charge is conserved, and the mass spectrum is exponential. This module works in the continuous setting for real numbers and defines the scale defect structure, its charge, stability, and mass formula.

\begin{definition}[\lean{Defect.ScaleDefect}]
A \emph{scale defect} of period $\Delta \in \mathbb{R}$ consists of:
\begin{itemize}
\item A continuously followed log-scale $\mathrm{lift} : \mathbb{R} \to \mathbb{R}$.
\item A topological charge $\mathrm{winding} : \mathbb{Z}$.
\item The quasiperiodicity condition: $\forall x, \quad \mathrm{lift}(x+1) = \mathrm{lift}(x) + (\mathrm{winding} : \mathbb{R}) \cdot \Delta$.
\end{itemize}
Going once around a loop ($x \mapsto x+1$) returns the scale shifted by $\mathrm{winding}$ periods.
\end{definition}

\begin{theorem}[\lean{Defect.ScaleDefect.winding\_unique}]
Assume $\Delta \neq 0$. If two defects have the same lift, they have the same winding: charge is well-defined.
\end{theorem}

\begin{proof}
Suppose $D$ and $E$ have the same lift. Evaluating the quasiperiodicity at $x = 0$:
\[ D.\mathrm{lift}(1) = D.\mathrm{lift}(0) + (D.\mathrm{winding} : \mathbb{R}) \cdot \Delta, \]
\[ E.\mathrm{lift}(1) = E.\mathrm{lift}(0) + (E.\mathrm{winding} : \mathbb{R}) \cdot \Delta. \]
Since the lifts are equal, $(D.\mathrm{winding} : \mathbb{R}) \cdot \Delta = (E.\mathrm{winding} : \mathbb{R}) \cdot \Delta$. Dividing by $\Delta \neq 0$ gives $(D.\mathrm{winding} : \mathbb{R}) = (E.\mathrm{winding} : \mathbb{R})$, so $D.\mathrm{winding} = E.\mathrm{winding}$.
\end{proof}

\begin{definition}[\lean{Defect.ScaleDefect.vacuum}]
The vacuum defect has zero lift and zero winding: it is the trivial configuration.
\end{definition}

\begin{theorem}[\lean{Defect.ScaleDefect.vacuum\_winding}]
The vacuum defect carries zero charge: $\mathrm{vacuum}.\mathrm{winding} = 0$.
\end{theorem}
\begin{proof}
Immediate from the definition of $\mathrm{vacuum}$, whose winding field is
set to $0$.
\end{proof}

\begin{definition}[\lean{Defect.ScaleDefect.combine}]
The superposition of two defects $D$ and $E$ is the defect with lift $D.\mathrm{lift} + E.\mathrm{lift}$ and winding $D.\mathrm{winding} + E.\mathrm{winding}$.
\end{definition}

\begin{theorem}[\lean{Defect.ScaleDefect.combine\_winding}]
Charge is additive: $\mathrm{combine}(D, E).\mathrm{winding} = D.\mathrm{winding} + E.\mathrm{winding}$.
\end{theorem}

\begin{proof}
Immediate from the definition of $\mathrm{combine}$.
\end{proof}

\begin{definition}[\lean{Defect.ScaleDefect.anti}]
The antiparticle of a defect $D$ is obtained by negating the lift and the winding.
\end{definition}

\begin{theorem}[\lean{Defect.ScaleDefect.anti\_winding}]
The antiparticle carries opposite charge: $D.\mathrm{anti}.\mathrm{winding} = -D.\mathrm{winding}$.
\end{theorem}

\begin{proof}
Immediate from the definition of $\mathrm{anti}$.
\end{proof}

\begin{theorem}[\lean{Defect.ScaleDefect.pair\_winding\_zero}]
A particle together with its antiparticle carries zero total charge: $\mathrm{combine}(D, D.\mathrm{anti}).\mathrm{winding} = 0$.
\end{theorem}

\begin{proof}
\[ \mathrm{combine}(D, D.\mathrm{anti}).\mathrm{winding} = D.\mathrm{winding} + D.\mathrm{anti}.\mathrm{winding} = D.\mathrm{winding} + (-D.\mathrm{winding}) = 0. \]
\end{proof}

\begin{definition}[\lean{Defect.ScaleDefect.IsTrivial}]
A defect is \emph{trivial} if the scale genuinely closes up around the loop: $\forall x, \quad D.\mathrm{lift}(x+1) = D.\mathrm{lift}(x)$.
\end{definition}

\begin{theorem}[\lean{Defect.ScaleDefect.vacuum\_isTrivial}]
The vacuum defect is trivial.
\end{theorem}

\begin{proof}
The vacuum lift is constant, so it trivially satisfies the periodicity condition.
\end{proof}

\begin{theorem}[\lean{Defect.ScaleDefect.stable\_of\_winding\_ne\_zero}]
Assume $\Delta \neq 0$. A defect with nonzero winding is not trivial; the scale genuinely fails to close, and no continuous deformation preserving the winding can make it close.
\end{theorem}

\begin{proof}
Suppose $D.\mathrm{winding} \neq 0$ but $D$ is trivial. Then $D.\mathrm{lift}(1) = D.\mathrm{lift}(0)$. But the quasiperiodicity gives $D.\mathrm{lift}(1) = D.\mathrm{lift}(0) + (D.\mathrm{winding} : \mathbb{R}) \cdot \Delta$, so $(D.\mathrm{winding} : \mathbb{R}) \cdot \Delta = 0$. Since $\Delta \neq 0$, we have $D.\mathrm{winding} = 0$, contradicting the assumption.
\end{proof}

\begin{theorem}[\lean{Defect.ScaleDefect.winding\_eq\_zero\_of\_trivial}]
Assume $\Delta \neq 0$. A trivial configuration has zero charge.
\end{theorem}

\begin{proof}
Follows by contrapositive from the previous theorem.
\end{proof}

\begin{definition}[\lean{Defect.ScaleDefect.mass}]
The mass of a defect with topological charge $k \in \mathbb{Z}$ is
\[ m(k) = m_0 \exp(|k| \cdot \Delta), \]
where $m_0$ is the mass of the vacuum. This is a posited law relating topology to dynamics; it is not derived from the axioms.
\end{definition}

\begin{theorem}[\lean{Defect.ScaleDefect.mass\_pos}]
If $m_0 > 0$, then the mass of any defect is positive.
\end{theorem}

\begin{proof}
The exponential is always positive, and the product of positive numbers is positive.
\end{proof}

\begin{theorem}[\lean{Defect.ScaleDefect.antiparticle\_same\_mass}]
A particle and its antiparticle have exactly equal mass: $m(-k) = m(k)$ for all $k \in \mathbb{Z}$.
\end{theorem}

\begin{proof}
\[ m(-k) = m_0 \exp(|-k| \cdot \Delta) = m_0 \exp(|k| \cdot \Delta) = m(k). \]
\end{proof}

\begin{theorem}[\lean{Defect.ScaleDefect.mass\_zero}]
The vacuum (charge zero) has the baseline mass $m_0$.
\end{theorem}

\begin{proof}
\[ m(0) = m_0 \exp(|0| \cdot \Delta) = m_0 \exp(0) = m_0. \]
\end{proof}

\begin{theorem}[\lean{Defect.ScaleDefect.mass\_ratio\_geometric}]
For $k \geq 0$, consecutive mass levels differ by the constant factor $e^\Delta$:
\[ m(k+1) = e^\Delta \cdot m(k). \]
The tower is geometric, given the posited mass law.
\end{theorem}

\begin{proof}
For $k \geq 0$, $|k+1| = k+1$ and $|k| = k$, so
\begin{align}
m(k+1) &= m_0 \exp((k+1) \cdot \Delta)\\
&= m_0 \exp(k \cdot \Delta + \Delta)\\
&= m_0 \exp(k \cdot \Delta) \cdot \exp(\Delta)\\
&= e^\Delta \cdot m(k).
\end{align}
\end{proof}

\begin{theorem}[\lean{Defect.ScaleDefect.mass\_ratio\_const}]
The ratio between consecutive levels is the same everywhere in the tower: for any $k, l \geq 0$ with $m_0 \neq 0$,
\[ \frac{m(k+1)}{m(k)} = \frac{m(l+1)}{m(l)}. \]
\end{theorem}

\begin{proof}
By the previous theorem, both ratios equal $e^\Delta$, independent of $k$ and $l$.
\end{proof}

\subsection{Charges}

Point defects in a uniaxial scale pattern carry two independent topological charges. The block charge lives in a finite group (the symmetries of the degenerate block) and is confined: a nontrivial block charge is not observable alone. The hedgehog charge is an integer, and no observability condition constrains it. This module establishes the two charges, their independence, and the asymmetry of confinement.

\begin{definition}[\lean{Charges.DefectCharge}]
A \emph{defect charge} in a uniaxial scale pattern is a pair $(b, q)$ where:
\begin{itemize}
\item $b \in B$ is the \emph{block charge}, a group element of the finite symmetry group $B$ of the degenerate directions.
\item $q \in \mathbb{Z}$ is the \emph{hedgehog charge}, the integer wrapping of the axis over the enclosing sphere.
\end{itemize}
\end{definition}

\begin{definition}[\lean{Charges.DefectCharge.comp}]
Two defects compose by multiplying their block charges and adding their hedgehog charges:
\[ (b, q) \circ (b', q') = (b \cdot b', q + q'). \]
\end{definition}

\begin{theorem}[\lean{Charges.DefectCharge.comp\_block}]
The block component of a composition is multiplicative: $\mathrm{comp}(c, d).\mathrm{block} = c.\mathrm{block} \cdot d.\mathrm{block}$.
\end{theorem}

\begin{proof}
Immediate from the definition of $\mathrm{comp}$.
\end{proof}

\begin{theorem}[\lean{Charges.DefectCharge.comp\_hedgehog}]
The hedgehog component of a composition is additive: $\mathrm{comp}(c, d).\mathrm{hedgehog} = c.\mathrm{hedgehog} + d.\mathrm{hedgehog}$.
\end{theorem}

\begin{proof}
Immediate from the definition of $\mathrm{comp}$.
\end{proof}

\begin{definition}[\lean{Charges.DefectCharge.anti}]
The antidefect of charge $(b, q)$ is $(b^{-1}, -q)$.
\end{definition}

\begin{theorem}[\lean{Charges.DefectCharge.anti\_block}]
The block component of an antidefect is inverse: $\mathrm{anti}(c).\mathrm{block} = c.\mathrm{block}^{-1}$.
\end{theorem}

\begin{proof}
Immediate from the definition of $\mathrm{anti}$.
\end{proof}

\begin{theorem}[\lean{Charges.DefectCharge.anti\_hedgehog}]
The hedgehog component of an antidefect is negated: $\mathrm{anti}(c).\mathrm{hedgehog} = -c.\mathrm{hedgehog}$.
\end{theorem}

\begin{proof}
Immediate from the definition of $\mathrm{anti}$.
\end{proof}

\begin{definition}[\lean{Charges.DefectCharge.triv}]
The vacuum has trivial block charge and zero hedgehog charge: $(1, 0)$.
\end{definition}

\begin{theorem}[\lean{Charges.DefectCharge.triv\_observable}]
The vacuum is observable.
\end{theorem}

\begin{proof}
Its block charge is the identity.
\end{proof}

\begin{definition}[\lean{Charges.DefectCharge.Observable}]
A defect charge $(b, q)$ is \emph{observable} if its block component is the identity: $b = 1$. There is no constraint on $q$.
\end{definition}

\begin{theorem}[\lean{Charges.DefectCharge.block\_confined}]
A defect with nontrivial block charge ($b \neq 1$) is not observable.
\end{theorem}

\begin{proof}
Observability requires $b = 1$ by definition.
\end{proof}

\begin{theorem}[\lean{Charges.DefectCharge.hedgehog\_free}]
For every integer $k \in \mathbb{Z}$, there exists an observable defect carrying that hedgehog charge: $\exists c, \, \mathrm{Observable}(c) \land c.\mathrm{hedgehog} = k$.
\end{theorem}

\begin{proof}
The charge $(1, k)$ is observable and has hedgehog charge $k$.
\end{proof}

\begin{theorem}[\lean{Charges.DefectCharge.hedgehog\_surjective\_on\_observables}]
Every integer hedgehog charge is realized by some observable state; the integral charge is completely unconstrained by observability.
\end{theorem}

\begin{proof}
This is a restatement of the previous theorem.
\end{proof}

\begin{theorem}[\lean{Charges.DefectCharge.exactly\_one\_confined}]
The two charges are asymmetric:
\begin{itemize}
\item Every observable defect has trivial block charge.
\item Every integer hedgehog charge appears in some observable defect.
\end{itemize}
Confinement applies only to the block charge.
\end{theorem}

\begin{proof}
Both statements follow from the definition of observability and the previous two theorems.
\end{proof}

\begin{theorem}[\lean{Charges.DefectCharge.hedgehog\_not\_determined\_by\_block}]
Fixing the block charge leaves the hedgehog charge free. For any group element $b$ and distinct integers $k, l$, there exist defects with the same block charge but different hedgehog charges.
\end{theorem}

\begin{proof}
The charges $(b, k)$ and $(b, l)$ have the same block component but different hedgehog components.
\end{proof}

\begin{theorem}[\lean{Charges.DefectCharge.block\_not\_determined\_by\_hedgehog}]
Fixing the hedgehog charge leaves the block charge free. For any integers $k$ and distinct group elements $b, b'$, there exist defects with the same hedgehog charge but different block charges.
\end{theorem}

\begin{proof}
The charges $(b, k)$ and $(b', k)$ have the same hedgehog component but different block components.
\end{proof}

\begin{theorem}[\lean{Charges.DefectCharge.charges\_independent}]
The two charges are genuinely independent: fixing one does not determine the other.
\end{theorem}

\begin{proof}
The previous two theorems establish this directly.
\end{proof}

\begin{theorem}[\lean{Charges.DefectCharge.pair\_observable}]
A particle and its antiparticle together carry trivial block charge, so the pair is observable.
\end{theorem}

\begin{proof}
If $c = (b, q)$, then $\mathrm{anti}(c) = (b^{-1}, -q)$, so $\mathrm{comp}(c, \mathrm{anti}(c)) = (b \cdot b^{-1}, q - q) = (1, 0)$, which is observable.
\end{proof}

\begin{theorem}[\lean{Charges.DefectCharge.pair\_hedgehog\_zero}]
A particle and its antiparticle together carry zero hedgehog charge.
\end{theorem}

\begin{proof}
If $c = (b, q)$, then $\mathrm{comp}(c, \mathrm{anti}(c)).\mathrm{hedgehog} = q + (-q) = 0$.
\end{proof}

\begin{theorem}[\lean{Charges.DefectCharge.observable\_pair\_any\_hedgehog}]
An observable pair of defects can carry any total hedgehog charge, provided their block charges cancel. If $c = (b, k)$ and $d = (b^{-1}, l)$, then $\mathrm{comp}(c, d)$ is observable with hedgehog charge $k + l$.
\end{theorem}

\begin{proof}
$\mathrm{comp}(c, d) = (b \cdot b^{-1}, k + l) = (1, k+l)$, which is observable.
\end{proof}

\subsection{Particle}

The algebra constrains the scale pattern dramatically: all `n` independent directions collapse to a single vector in the scaling sector. The rank-one structure gives one exponential rate, and rank two or higher necessarily produce mixtures with $\mathrm{CV}^2 > 1$. The observed spectrum has $\mathrm{CV}^2 < 1$, disfavouring higher rank. The order parameter is a projective axis, so the block group is $\mathbb{Z}/2$, giving a two-valued confined charge. A particle is a triple: threshold, $\mathbb{Z}/2$ label, and integer charge.

\begin{theorem}[\lean{Particle.boost\_injective}]
Scalings are faithfully labelled by their direction vector: if $\mathrm{boost}(v) = \mathrm{boost}(w)$ for $v, w \in \text{Fin } k \to \mathbb{R}$, then $v = w$.
\end{theorem}

\begin{proof}
If $\mathrm{boost}(v) = \mathrm{boost}(w)$, then the boost elements agree. The boost at direction index $(0, i+1)$ recovers $v_i$, so $v_i = w_i$ for all $i$.
\end{proof}

\begin{theorem}[\lean{Particle.pattern\_is\_one\_vector}]
The scale pattern is one vector in the scaling sector, and scaling it is the only freedom: $\mathrm{boost}(c \cdot v) = c \cdot \mathrm{boost}(v)$ for all $c \in \mathbb{R}$ and $v \in \text{Fin } k \to \mathbb{R}$.
\end{theorem}

\begin{proof}
Follows from the linearity of the boost construction.
\end{proof}

\begin{theorem}[\lean{Particle.mixture\_second\_moment}]
For a two-component mixture with weights $p, 1-p$ and means $a, b$, the excess of the second moment over the squared mean is
\[ p a^2 + (1-p)b^2 - (pa + (1-p)b)^2 = p(1-p)(a-b)^2. \]
\end{theorem}

\begin{proof}
Expand both sides:
\begin{align}
\text{LHS} &= pa^2 + (1-p)b^2 - p^2a^2 - 2p(1-p)ab - (1-p)^2b^2\\
&= pa^2(1 - p) + b^2(1-p)(1 - (1-p)) - 2p(1-p)ab\\
&= p(1-p)a^2 + p(1-p)b^2 - 2p(1-p)ab\\
&= p(1-p)(a^2 + b^2 - 2ab)\\
&= p(1-p)(a-b)^2.
\end{align}
\end{proof}

\begin{definition}[\lean{Particle.mixCvSq}]
The squared coefficient of variation of a mixture of two exponential components with means $a, b$ and weights $p, 1-p$ is
\[ \mathrm{CV}^2_{\mathrm{mix}}(p, a, b) = \frac{2(pa^2 + (1-p)b^2)}{(pa + (1-p)b)^2} - 1. \]
An exponential distribution has second moment twice its squared mean, so this formula gives the mixture's squared coefficient of variation.
\end{definition}

\begin{theorem}[\lean{Particle.mixture\_cvSq\_ge\_one}]
Any mixture of two exponential rates with weights in $[0,1]$ and positive weighted mean has $\mathrm{CV}^2_{\mathrm{mix}} \geq 1$.
\end{theorem}

\begin{proof}
By the identity above, the numerator $2(pa^2 + (1-p)b^2) - (pa + (1-p)b)^2$ contains the term $2p(1-p)(a-b)^2 \geq 0$. The denominator is positive, so
\[ \mathrm{CV}^2_{\mathrm{mix}} = \frac{2(pa^2 + (1-p)b^2)}{(pa+(1-p)b)^2} - 1 \geq 1. \]
\end{proof}

\begin{theorem}[\lean{Particle.mixture\_cvSq\_eq\_one\_iff}]
For a proper mixture ($0 < p < 1$) with positive weighted mean, $\mathrm{CV}^2_{\mathrm{mix}} = 1$ if and only if $a = b$ (the two rates coincide).
\end{theorem}

\begin{proof}
If $a = b$, then $(a-b)^2 = 0$, so the numerator becomes $2a^2$ and the denominator becomes $a^2$, giving $\mathrm{CV}^2_{\mathrm{mix}} = 2 - 1 = 1$.

Conversely, if $\mathrm{CV}^2_{\mathrm{mix}} = 1$, then from the mixture identity, the term $p(1-p)(a-b)^2$ must vanish. Since $0 < p < 1$, we have $p(1-p) > 0$, so $(a-b)^2 = 0$ and $a = b$.
\end{proof}

\begin{theorem}[\lean{Particle.observed\_disfavours\_higher\_rank}]
The observed spectrum has $\mathrm{CV}^2 < 0.88$. Since every mixture of two or more rates satisfies $\mathrm{CV}^2 \geq 1$, the observed spectrum is evidence against a higher-rank algebra and in favour of rank one.
\end{theorem}

\begin{proof}
Let $\mathrm{CV}^2_{\mathrm{obs}} < 0.88$ and $\mathrm{CV}^2_{\mathrm{mix}} \geq 1$ for any mixture. Then $\mathrm{CV}^2_{\mathrm{obs}} < \mathrm{CV}^2_{\mathrm{mix}}$.
\end{proof}

\begin{theorem}[\lean{Particle.rate\_value\_not\_fixed}]
The value of the exponential rate is not determined by the framework. The coefficient of variation is invariant under rescaling both component means: $\mathrm{CV}^2_{\mathrm{mix}}(p, ca, cb) = \mathrm{CV}^2_{\mathrm{mix}}(p, a, b)$ for all $c \neq 0$.
\end{theorem}

\begin{proof}
Scaling both means by $c$ scales both the numerator and denominator by $c^2$, leaving the ratio unchanged.
\end{proof}

\begin{definition}[\lean{Particle.Block}]
The block group of the uniaxial pattern is the multiplicative group on $\mathbb{Z}/2$, denoted $\mathrm{Multiplicative}(\mathbb{Z}/2)$.
\end{definition}

\begin{definition}[\lean{Particle.Charge}]
A \emph{particle charge} is a defect charge in the block group $\mathrm{Multiplicative}(\mathbb{Z}/2)$.
\end{definition}

\begin{theorem}[\lean{Particle.block\_two\_valued}]
The block group has cardinality 2.
\end{theorem}

\begin{proof}
$\mathbb{Z}/2$ has two elements, and the multiplicative group on it has the same cardinality.
\end{proof}

\begin{theorem}[\lean{Particle.block\_cannot\_be\_three\_valued}]
The framework structurally cannot produce a three-valued confined charge: the block group is $\mathbb{Z}/2$, which has exactly 2 elements, not 3.
\end{theorem}

\begin{proof}
The order parameter is projective (unoriented axis), and the loop homotopy group of projective space is $\mathbb{Z}/2$. This is a topological fact, not a free parameter.
\end{proof}

\begin{theorem}[\lean{Particle.block\_self\_inverse}]
Every element of the block group is its own inverse: $b \cdot b = 1$ for all $b \in \mathrm{Block}$.
\end{theorem}

\begin{proof}
In $\mathbb{Z}/2$, every nonidentity element satisfies $b^2 = 1$.
\end{proof}

\begin{definition}[\lean{Particle.Species}]
A \emph{particle species} is determined by:
\begin{itemize}
\item A threshold $t \in \mathbb{R}$, the log-scale at which the particle becomes resolvable.
\item A charge $(b, q) \in \mathrm{Charge}$, specifying the $\mathbb{Z}/2$ label and the integer charge.
\end{itemize}
The framework supplies no further labels.
\end{definition}

\begin{definition}[\lean{Particle.bind}]
Two particles with charges $c_1, c_2$ can be bound at a new threshold $t$ to form a composite particle with the same threshold and charge $\mathrm{comp}(c_1, c_2)$.
\end{definition}

\begin{theorem}[\lean{Particle.bind\_hedgehog}]
The hedgehog charge of a composite is the sum of the components' hedgehog charges.
\end{theorem}

\begin{proof}
By definition of $\mathrm{comp}$.
\end{proof}

\begin{theorem}[\lean{Particle.bind\_block}]
The block charge of a composite is the product of the components' block charges.
\end{theorem}

\begin{proof}
By definition of $\mathrm{comp}$.
\end{proof}

\begin{theorem}[\lean{Particle.two\_odd\_make\_even}]
Two $\mathbb{Z}/2$-odd particles (same block charge) compose to an even particle (trivial block charge), without needing any dynamics: $\mathrm{bind}(x, y, t).\mathrm{charge}.\mathrm{block} = 1$ whenever $x.\mathrm{charge}.\mathrm{block} = y.\mathrm{charge}.\mathrm{block}$.
\end{theorem}

\begin{proof}
If both have block charge $b$, the composite has block charge $b \cdot b = 1$ by \lean{Particle.block\_self\_inverse}.
\end{proof}

\begin{theorem}[\lean{Particle.hedgehog\_conserved\_in\_splitting}]
The hedgehog charge is exactly conserved in any splitting: if a particle splits into two, the parts' hedgehog charges sum to the whole's. This is an identity, independent of dynamics.
\end{theorem}

\begin{proof}
If the whole has charge $(b, q)$ and splits into $(b_1, q_1)$ and $(b_2, q_2)$ with $b = b_1 \cdot b_2$ and $q = q_1 + q_2$, then $q - q_1 = q_2$, which is identically true.
\end{proof}

\begin{theorem}[\lean{Particle.no\_further\_label}]
A particle is determined by its threshold and its two charges: two particles agreeing on all three are equal.
\end{theorem}

\begin{proof}
The species structure has two fields (threshold and charge), and the charge has two components (block and hedgehog). Two species with equal thresholds and charges are equal by extensionality.
\end{proof}

\begin{theorem}[\lean{Particle.particle\_label\_structure}]
Every particle has the form $(t, b, q)$ for some threshold $t \in \mathbb{R}$, block label $b \in \mathbb{Z}/2$, and integer charge $q \in \mathbb{Z}$.
\end{theorem}

\begin{proof}
The species and charge structures decompose as stated.
\end{proof}

\subsection{Emergence}

Particles are not fundamental: they are features of a slice across the scale axis. By axiom A3, energy and mass are inverse scale; so the set of observable particles depends on the resolution at which the universe is probed. The framework makes no claim about what particles exist, only that at each scale a new content emerges, growing monotonically without bound. This module characterizes the threshold structure and establishes why there is no complete list of ``fundamental particles.''

\begin{definition}[\lean{Emergence.resolved}]
The set of particles resolvable at log-scale $t$ is $\{k \mid \mu_k \leq t\}$, where $\mu$ is the threshold sequence.
\end{definition}

\begin{theorem}[\lean{Emergence.mem\_resolved}]
An index $k$ is in the resolved set at scale $t$ if and only if its threshold satisfies $\mu_k \leq t$.
\end{theorem}

\begin{proof}
Immediate from the definition.
\end{proof}

\begin{theorem}[\lean{Emergence.resolved\_mono}]
Content only grows with better resolution: if $t \leq t'$, then $\mathrm{resolved}(\mu, t) \subseteq \mathrm{resolved}(\mu, t')$.
\end{theorem}

\begin{proof}
If $k \in \mathrm{resolved}(\mu, t)$, then $\mu_k \leq t \leq t'$, so $k \in \mathrm{resolved}(\mu, t')$.
\end{proof}

\begin{theorem}[\lean{Emergence.content\_never\_complete}]
If the thresholds are unbounded (for every upper bound $M$ there exists a threshold above it), then no scale sees all particles: $\mathrm{resolved}(\mu, t) \neq \text{all indices}$ for any $t$.
\end{theorem}

\begin{proof}
Given $t$, there exists $k$ with $t < \mu_k$, so $k \notin \mathrm{resolved}(\mu, t)$.
\end{proof}

\begin{theorem}[\lean{Emergence.always\_more\_above}]
The tower never terminates: for any scale $t$, there exists a scale $t' > t$ at which strictly more particles are resolved.
\end{theorem}

\begin{proof}
Given $t$, let $k$ satisfy $t < \mu_k$ (which exists by unboundedness). Then $t < \mu_k \leq t'$ for $t' = \mu_k$ shows that $k \in \mathrm{resolved}(\mu, t')$ but $k \notin \mathrm{resolved}(\mu, t)$.
\end{proof}

\begin{theorem}[\lean{Emergence.no\_fundamental\_list}]
The content is a never-completed, strictly increasing family: for all scales $t$, the resolved content differs from the full set, and there always exists a higher scale with strictly larger content.
\end{theorem}

\begin{proof}
Follows from the previous two theorems. The question ``which particles are fundamental?'' has no scale-independent answer; it is a malformed question in this framework.
\end{proof}

\begin{definition}[\lean{Emergence.massOf}]
The mass of a particle with threshold $\mu_k$ is $m_k = e^{\mu_k}$. By axiom A3, energy and mass are inverse scale, so higher threshold means heavier particle.
\end{definition}

\begin{theorem}[\lean{Emergence.mass\_iff\_threshold}]
Mass and threshold determine each other: $\mathrm{massOf}(\mu, k) = \mathrm{massOf}(\mu, l) \iff \mu_k = \mu_l$.
\end{theorem}

\begin{proof}
The exponential function is injective: $e^x = e^y \iff x = y$.
\end{proof}

\begin{theorem}[\lean{Emergence.mass\_mono}]
Heavier particles appear at higher thresholds: if $\mu_k < \mu_l$, then $\mathrm{massOf}(\mu, k) < \mathrm{massOf}(\mu, l)$.
\end{theorem}

\begin{proof}
The exponential function is strictly increasing.
\end{proof}

\begin{theorem}[\lean{Emergence.content\_eq\_of\_no\_threshold}]
If there is no threshold in the interval $(t, t']$, then the resolved content at $t$ and $t'$ is identical: the effective description is exact between thresholds.
\end{theorem}

\begin{proof}
For any index $k$, if $\mu_k \leq t$, then $\mu_k \leq t'$, so $k$ remains resolved. If $\mu_k > t'$, then by the no-threshold assumption, $\mu_k$ is not in $(t, t']$, so $\mu_k > t'$ and $k$ is not resolved at either scale. No index crosses the gap.
\end{proof}

\begin{theorem}[\lean{Emergence.content\_changes\_only\_at\_threshold}]
Content changes between two scales if and only if at least one threshold lies in the gap: if $\mathrm{resolved}(\mu, t) \neq \mathrm{resolved}(\mu, t')$ for $t \leq t'$, then there exists $k$ with $t < \mu_k \leq t'$.
\end{theorem}

\begin{proof}
If no such $k$ exists, then by the previous theorem the contents are equal, giving a contradiction. Therefore at least one threshold lies in the gap.
\end{proof}

